\subsection{Collective action, sanctions, and reputation}

Collective-action problems arise when private incentives erode shared returns \pcite{olson1965logic,hardin1968tragedy}.
Taking an empirical lens, laboratory and evolutionary traditions show that monitoring, costly punishment, and reputation can sustain cooperation \pcite{fehr2002altruistic,boyd2003evolution,axelrod1981evolution,axelrod1984evolution}.
More precisely, indirect reciprocity and image scoring formalize how observed giving shapes future partners \pcite{wedekind2000cooperation,panchanathan2004indirect}.
Looking at this we can see that reputation helps solve commons dilemmas by linking contribution to standing and is among the core evolutionary routes to cooperation \pcite{milinski2002reputation,nowak2006five}.
This is important because gossip can transmit reputation beyond direct observation \pcite{sommerfeld2007gossip}.
Across these lines, the working hypothesis is directional: more social visibility should raise, or at least protect, cooperative transfers.

Beyond informal sanctions, institutional design complements reputation and Ostrom documents polycentric rules, graduated sanctions, and participatory governance in durable commons regimes \pcite{ostrom1990governing}.
Those design principles hence motivate dual channels in computational societies: formal enforcement inside one institution and social evaluation everywhere, with periodic voting on fiscal parameters.
Yet we observe that classical literatures rarely ask what happens when social marks land on agents who cannot exit and cannot punish peers.

\subsection{Language-model agents as strategic laboratories}

Recent work operationalizes these dynamics in language-model agents, which produce numeric choices together with natural-language rationales, belief updates, and peer judgments \pcite{park2023generative,sumers2024cognitive}.
Consequently, game-theoretic LLM platforms treat such agents as laboratories for strategic interaction while cooperation benchmarks report heterogeneous success across games, prompts, and communication regimes \pcite{mao2023alympics, piatti2024govsim,madmoun2025communication,liang2025everyone}.
On balance some societies sustain contributions when talk is rich; others collapse once incentives dominate.

Interestingly enough, a growing strand shows that explicit reasoning can reduce cooperation.
Specifically, models that deliberate more intensely often favor individually rational free-riding, mirroring human ``spontaneous giving and calculated greed'' \pcite{li2025spontaneous,freeriders2025,rand2012spontaneous}.
Relatedly, norm-focused multi-agent studies report partial emergence and fragile stability rather than automatic internalization \pcite{gupta2025norms,ashery2025conventions,cross2025norms}.
Likewise, human public-goods experiments show conditional cooperation and intuition--deliberation tensions that foreshadow these LLM patterns \pcite{fischbacher2001conditional,rand2012spontaneous}.
All in all, these results motivate separating average transfer levels from norm strength.

\subsection{Reputation as an assumed stabilizer}

Returning to reputation specifically, in several recent LLM-agent studies, reputation or peer talk are treated as stabilizers of cooperation. \pcite{repunet2025,piatti2024govsim,madmoun2025communication}.
Taking from classical evolutionary models, the punishment-strategy transplants  show that explicit sanctions can organize behavior when costs and targets are clear \pcite{warnakulasuriya2025punishment}.
However, where gossip substitutes for formal punishment, the mapping from stigma to repair is unsettled.
Compounding this uncertainty, self-reinforcing salience loops in LLM inference can lock opportunistic frames once they appear \pcite{howlround2025}.
Prior LLM work rarely event-studies contribution intensity after negative social marks under forced institutional partitions.

\subsection{Climate risk sharing and the research gap}

Climate governance couples mitigation public goods with asymmetric vulnerability and compensation claims.
Loss-and-damage instruments are dual-use in spirit: contributions can fund shared returns and contingent payouts after shocks.
Computational societies that combine those fiscal rules with LLM agents that gossip and vote remain rare.
Endogenous democracy lets agents revise punishment, subsidy, and equity weights, which is politically meaningful and confounds clean shock identification when votes coincide with damage rounds.

Consequently, prior work largely assumes that reputation and gossip raise cooperation, or at least do not systematically reduce it.
LLM-agent papers that celebrate communication rarely measure post-stigma intensity under climate asymmetry across independent replications.
We address that gap with an event study of contribution intensity after reputation and gossip events in a dual-replication climate-risk society.
